% Options for packages loaded elsewhere
\PassOptionsToPackage{unicode}{hyperref}
\PassOptionsToPackage{hyphens}{url}
%
\documentclass[
  10pt,
  a4paper,
]{article}
\usepackage{amsmath,amssymb}
\usepackage{iftex}
\ifPDFTeX
  \usepackage[T1]{fontenc}
  \usepackage[utf8]{inputenc}
  \usepackage{textcomp} % provide euro and other symbols
\else % if luatex or xetex
  \usepackage{unicode-math} % this also loads fontspec
  \defaultfontfeatures{Scale=MatchLowercase}
  \defaultfontfeatures[\rmfamily]{Ligatures=TeX,Scale=1}
\fi
\IfFileExists{lmodern.sty}{\usepackage{lmodern}}{}
\ifPDFTeX\else
  % xetex/luatex font selection
\fi
% Use upquote if available, for straight quotes in verbatim environments
\IfFileExists{upquote.sty}{\usepackage{upquote}}{}
\IfFileExists{microtype.sty}{% use microtype if available
  \usepackage[]{microtype}
  \UseMicrotypeSet[protrusion]{basicmath} % disable protrusion for tt fonts
}{}
\makeatletter
\@ifundefined{KOMAClassName}{% if non-KOMA class
  \IfFileExists{parskip.sty}{%
    \usepackage{parskip}
  }{% else
    \setlength{\parindent}{0pt}
    \setlength{\parskip}{6pt plus 2pt minus 1pt}}
}{% if KOMA class
  \KOMAoptions{parskip=half}}
\makeatother
\setlength{\parskip}{4pt plus 1pt minus 1pt}
\usepackage{xcolor}
\usepackage[margin=2.2cm]{geometry}
\usepackage{longtable,booktabs,array}
\usepackage{calc} % for calculating minipage widths
% Correct order of tables after \paragraph or \subparagraph
\usepackage{etoolbox}
\makeatletter
\patchcmd\longtable{\par}{\if@noskipsec\mbox{}\fi\par}{}{}
\makeatother
% Allow footnotes in longtable head/foot
\IfFileExists{footnotehyper.sty}{\usepackage{footnotehyper}}{\usepackage{footnote}}
\makesavenoteenv{longtable}
\setlength{\emergencystretch}{3em} % prevent overfull lines
\providecommand{\tightlist}{%
  \setlength{\itemsep}{0pt}\setlength{\parskip}{0pt}}
\setcounter{secnumdepth}{-\maxdimen} % remove section numbering
\ifLuaTeX
  \usepackage{selnolig}  % disable illegal ligatures
\fi
\usepackage{bookmark}
\IfFileExists{xurl.sty}{\usepackage{xurl}}{} % add URL line breaks if available
\urlstyle{same}
\hypersetup{
  hidelinks,
  pdfcreator={LaTeX via pandoc}}

\author{}
\date{}

\begin{document}

\section{Response to Reviewers ---
TCS-D-26-00494}\label{response-to-reviewers-tcs-d-26-00494}

Manuscript: ``Distributional Learning of Context-Free Languages under
Fixed Finite-Monoid Typing''

Author: Takayuki Kuriyama

We thank the Editor and both reviewers for their detailed and constructive reports. Their comments led us to correct the learning algorithm, remove an unused typing component, repair the terminal-rule completeness argument, and substantially reorganize the manuscript.

Throughout this response, unqualified section, theorem, proposition, lemma,
and corollary numbers refer to the revised manuscript. References to the
submitted manuscript are explicitly marked ``submitted'' or ``former'';
reviewer quotations retain their original numbering. Quoted excerpts reproduce
the reviewers' wording; list formatting and LaTeX notation may be normalized
typographically, and omissions are indicated by \ldots{}.  Each reviewer
comment is numbered below for traceability.  The structural revision locators
(section/theorem/lemma) are supplemented by page and printed-line anchors from
the author-compiled v82 manuscript PDF (26 pages).  If pagination changes
during the final submission-format pass, these mechanical page/line anchors
will be regenerated before upload.

\subsection{Overview of changes}\label{overview-of-changes}

The revision makes eight substantive changes in response to the Editor
and reviewers. The first three correct the learner or its proof architecture;
the fourth clarifies the empty-word convention; the fifth isolates a quantitative
boundary for arbitrary fixed typings; the sixth makes the fixed-window
correspondence exact and verifies a fixed-window bound polynomial in the two
parameters used by Yoshinaka---grammar size and ordinary thickness; the seventh clarifies the genuine scope of
the linear theorem; and the eighth states precisely which polynomial guarantees
remain after separating the batch constructor from the sequential learner.

Three claims from the submitted version have been withdrawn or narrowed:
(i) order-independence of the sequential learner, (ii) the claimed necessity of
outer-context typing, and (iii) the ``polynomial time-and-data'' label for the
sequential linear learner. The revised manuscript retains only the claims
proved by the corrected architecture.

\begin{enumerate}
\def\labelenumi{\arabic{enumi}.}
\tightlist
\item
  \textbf{Learner correctness.} The submitted sequential learner did not
  necessarily stabilize syntactically. It has been replaced by a
  conservative Gold learner, formally separated from the finite-sample
  batch constructor.
\item
  \textbf{Outer-context typing.} The submitted outer-context type
  annotations were not consulted by the proof. They have been removed,
  leaving only the yield type actually required by completeness.
\item
  \textbf{Terminal-rule completeness.} The submitted terminal-rule
  argument was incomplete when a typed nonterminal carried several terminal
  rules. The witness set now contains $u_Xav_X$ for every typed terminal rule
  $X\to a$, and the proof uses equal-yield-type Rule (R3) before lexical
  Rule (R4) when the canonical yield differs from $a$.
\item
  \textbf{Nonempty-fragment convention.} Fixed-\(h\) substitutability is now
  stated only for nonempty internal fragments, with \(\lambda\) handled
  separately at the start symbol.
\item
  \textbf{General fixed-\(h\) quantitative boundary.} New Section 7 first
  isolates the typed-thickness parameter governing the canonical
  full-refinement witnesses. Proposition 7.3 gives a fixed two-element typing
  for which typed thickness grows exponentially relative to source grammar
  size and ordinary thickness. Theorem 7.7 then gives a stronger,
  operator-independent obstruction inside the set-driven characteristic-sample
  framework: for a different homomorphism into the same abstract two-element
  zero monoid, nested finite targets
  \(\mathcal T_n^-\subsetneq\mathcal T_n\) have reduced CFG presentations
  of polynomial size and ordinary thickness, yet every characteristic sample for
  \(\mathcal T_n\), under any set-driven operator admitting characteristic
  samples for both targets, must contain an exponentially long positive word.
  Consequently, no set-driven characteristic-sample scheme covering all of
  \(\mathcal C_{h_\circ}^{\mathrm{cf}}\) admits a polynomial bound in
  reduced CFG size and ordinary thickness. The fixed-window and linear
  settings retain the positive bounds stated later.
\item
  \textbf{Fixed-window correspondence and compatibility.} Proposition 3.3
  shows that, for every fixed \((k,\ell)\), the finite-monoid typing
  \(h_{k,\ell}\) captures exactly Yoshinaka's
  \((k,\ell)\)-substitutability condition, including \((0,0)\). Section 8
  verifies that the present batch reconstruction has characteristic data
  polynomial in the same two parameters used in Yoshinaka's analysis---grammar
  size and ordinary thickness---without claiming the same polynomial or normal
  form.
\item
  \textbf{Scope of the linear theorem.} Section 9 adds a nonregular
  linear language \(L_{\pm,e}\) that is fixed-\(h\) substitutable but is
  neither ordinarily substitutable nor \((k,\ell)\)-substitutable for any
  fixed \(k,\ell\).
\item
  \textbf{Polynomial guarantees.} We retain polynomial-time finite-sample
  reconstruction and the quantitative characteristic-data bounds proved for
  the fixed-window and linear cases, but we no longer attach the submitted
  Theorem 7.9's ``polynomial time-and-data'' label to the corrected
  conservative sequential learner itself. The latter is stated separately as
  a Gold learner whose reconstructions and updates are polynomial in the
  accumulated sample size.
\end{enumerate}

After these corrections, the presentation was substantially compressed
and reorganized. The main structural changes are summarized below.

\begin{longtable}[]{@{}
  >{\raggedright\arraybackslash}p{(\columnwidth - 2\tabcolsep) * \real{0.5000}}
  >{\raggedright\arraybackslash}p{(\columnwidth - 2\tabcolsep) * \real{0.5000}}@{}}
\toprule\noalign{}
\begin{minipage}[b]{\linewidth}\raggedright
Submitted version
\end{minipage} & \begin{minipage}[b]{\linewidth}\raggedright
Revised version
\end{minipage} \\
\midrule\noalign{}
\endhead
\bottomrule\noalign{}
\endlastfoot
Section 5 recomputed a growing grammar after every positive datum. &
Section 4 now presents the finite-sample batch constructor and the
conservative sequential Gold learner before any target refinement is
introduced; Section 5 introduces yield typing only for the completeness
simulation. This follows the constructor/learner distinction used in the
classical substitutability literature and adopts Reviewer 1's suggested
order of presentation. \\
Section 4 used yield type plus outer-context types. & Outer-context
annotations are removed; only the yield type actually consulted by the
completeness proof remains. \\
The submitted manuscript did not separate ordinary source thickness from the
thickness created by full yield typing. & New Section 7 introduces typed
thickness for the general fixed-\(h\) analysis and first gives an
exponential source-to-typed gap for one fixed two-element monoid. It then
proves a stronger set-driven lower bound using nested finite targets under
another fixed two-element typing: every characteristic sample for the larger
target must contain an exponentially long word for any set-driven operator
that has characteristic samples for both targets. The level-coded construction
now includes an explicit unique-parsing and replacement-site argument before
the substitutability lemma. Thus the limitation is not an artifact of the
canonical full-refinement witnesses or of the specific reconstruction
operator. \\
The submitted fixed-window comparison was only expressive. & Proposition
3.3 now identifies the fixed \((k,\ell)\) class exactly with the class induced
by \(h_{k,\ell}\). Revised Section 8 is framed as a compatibility result: it checks
that the batch reconstruction is polynomial in the same two parameters used
by Yoshinaka---grammar size and ordinary thickness---while allowing a different
normal form and polynomial degree. The two short witness-bound
proofs and the general thickness-preserving normalization are collected in
Appendix A. \\
The linear part contained a long normalization development. &
Proposition 9.1 and Lemma 9.2 remain in Section 9, while the detailed
linear normalization and short-yield proof are moved to Appendices B and C.
Lemma 9.2 directly bounds typed thickness, and the general typed-thickness
corollary then supplies the witness bound. Appendix A contains the
fixed-window typed-yield bound together with the general
thickness-preserving SSBNF normalization. \\
The submitted boundary material was spread across Sections 8--10 and two
appendices. & It is consolidated into Section 10; the submitted
closure/counterexample and exceptional CCL\(_1\) appendices are removed. Only a compact proposition for
\(\Delta:=\{a^n b^n:n\ge0\}\) and its Kleene closure \(\Delta^*\) is
retained. It shows that the general fixed-\(h\) theorem extends beyond its
linear subclass: for the explicitly defined $h_\star$,
\(\Delta^*\in\mathcal C_{h_\star}^{\mathrm{cf}}\setminus\mathsf{KL}\). \\
The capped-counter notation used CCL\(_p\) and the alphabet
\(\{d,u,;\}\). & It is renamed CTR\(_\rho\) over the clearer alphabet
\(\{\uparrow,\downarrow,\#\}\). \\
\end{longtable}

For three suggestions, a short statement remains in the main text because it
is used immediately by a later result.
Proposition 9.1 and Lemma 9.2 remain stated in the main text because
Theorem 9.3 uses their single-spine normal form and short-yield bound,
although their detailed proofs are now in Appendices B and C. The former
Lemma 9.11, reduced to the right-quotient statement actually used, is now
Lemma 10.7 and remains immediately before that sole use rather than being moved
to an appendix. Finally, the boundary discussion retains only the compact proposition for
\(\Delta^*\), which records both a nonlinear fixed-\(h\) example and the
separation \(\Delta^*\notin\mathsf{KL}\).

The revised manuscript separates finite-sample guarantees for the batch
reconstruction from Gold stabilization for the conservative sequential learner.
A compact doubling-grammar example explains why no grammar-size-only polynomial
characteristic-data bound is claimed for arbitrary CFG representations; the
revision now states explicitly that the singleton languages in this example are
fixed-$h$ substitutable for every $h$.

\subsection{Response to the Editor}\label{response-to-the-editor}

\textbf{E.1. Editor (quoted excerpt):} ``The reviewers agree that the
scientific contribution has merit. However, they both point to many
issues with the presentation, many of which leak also into the
scientific part of the paper and are not merely language revisions. This
is especially relevant in the learning part of the paper. \ldots{}
Specifically, any correctness issues should be carefully considered.''

\textbf{Response:} We agree. The submitted learner has been replaced by
a conservative Gold learner; the terminal-rule gap has been repaired in both
the witness set and the completeness proof; and the unused outer-context
annotations have been removed. Fixed-\(h\) substitutability is now stated for
nonempty internal fragments, with \(\lambda\) handled separately. Section 4
defines the batch operator and sequential learner before Section 5 introduces
target-side yield typing for completeness. Polynomial characteristic data are
attributed to the set-driven batch operator, while the sequential learner is
asserted to have Gold stabilization with polynomial-time updates.

\textbf{Revision location:} Definition 2.1; Section 4.2; Section 5.1 and Theorem 5.3; Corollary 5.5; Sections 7--9.
\emph{v82 PDF anchors: p.~4, ll.~139--146; pp.~7--8, ll.~280--286; pp.~8--9, ll.~289--330; p.~9, ll.~332--351; p.~10, ll.~361--375; pp.~10--18, ll.~394--654.}

\subsection{Reviewer 1}\label{reviewer-1}

\subsubsection{General assessment and learning
algorithm}\label{general-assessment-and-learning-algorithm}

\textbf{R1.1. Comment (quoted excerpt):} ``The main result of this submission
is an elegant generalization of the results of Clark and Eyraud (2007)
on the learning of substitutable context-free languages (SCFLs) and of
Yoshinaka (2008) on the learning of k,l-SCFLs. \ldots{} The only serious
error is in the definition of their learning algorithm, but this is
easily repairable by following prior work.''

\textbf{Response:} The reviewer is correct; this exposed a genuine
correctness problem in the submitted learner. The submitted manuscript
recomputed a grammar from all data at every stage; this need not
stabilize syntactically. The revision distinguishes a set-driven batch
operator \(\mathcal B_h(K)=\widehat G(K)\) from a
conservative sequential learner \(\mathcal A_h\). We now also cite
Eyraud--Heinz--Yoshinaka (2016), whose Appendix explicitly discusses why a
set-based characteristic-sample result need not transfer directly to an
incremental learner. The wrapper used here is the same conservative device
used by Clark--Eyraud's SGL (Algorithm 2) and Yoshinaka's
\((k,\ell)\)-SGL (Algorithm 1): a new hypothesis is constructed only when the next positive datum is not generated by the current one. The learner keeps its
current grammar whenever the new
positive datum is already generated and rebuilds only when the datum
witnesses that the current hypothesis is too small. Once the finite
witness set \(\mathsf W(\widetilde G)\) has appeared, the next rebuild,
if any, reconstructs the target language exactly. After that no positive
datum can trigger another change. Corollary 5.5 now states this without temporal ambiguity: among stages $n\ge n_0$ after the witness set is present, at most one hypothesis change can occur. If no rebuild occurs after the witness
set has appeared, the current hypothesis still contains every positive
datum seen so far: rebuild hypotheses are sample-consistent, and a
non-rebuild step occurs only when the newly presented datum is already
generated. Hence any target word missing from the current hypothesis has
not yet occurred in the text; because every target word eventually
appears, such a word would later force a rebuild. Therefore the current
hypothesis must already be exact. Thus the hypothesis grammar itself
eventually stabilizes. The submitted order-independence claim has also been removed: only the
batch operator is set-driven. We retain the precise polynomial guarantees proved
for the batch operator: reconstruction is polynomial in the finite sample size;
fixed-window targets have characteristic data polynomial in grammar size and ordinary thickness; and linear targets have characteristic data polynomial in a chosen linear
CFG representation. We do not attach the submitted Theorem 7.9's
``polynomial time-and-data'' label to the corrected sequential learner itself,
because the latter is a separate conservative Gold wrapper and the batch output
need not be a linear representation.

\textbf{Revision location:} Section 2.1; Section 4.2; Corollary 5.5; Theorem 8.3; Theorem 9.3 and the following clarification.
\emph{v82 PDF anchors: pp.~3--4, ll.~117--134; pp.~7--8, ll.~280--286; p.~10, ll.~361--375; p.~16, ll.~579--586; p.~17, ll.~613--622.}

\textbf{R1.2. Comment (quoted excerpt):} ``On the other hand, the presentation leaves
considerable room for improvement. By addressing the absence of definitions for
technical terms, unnecessarily long proofs, inconsistent use of metavariables, and
the over-fragmentation of the paper into too many sections, the authors could
substantially reduce the burden on the reader, and I recommend a major revision
along these lines. \ldots{} I suspect that this material could be presented in less
than half the current length.''

\textbf{Response:} Agreed. The manuscript has been substantially shortened and reorganized. Definitions and the learning model are consolidated in Section 2, proof-local terminology has been removed or introduced only at the point of use, and avoidable metavariable collisions have been removed. The quantitative sections now reuse one typed-thickness engine rather than recounting witness sizes, the new lower bound is presented with compact reduced CFGs instead of auxiliary SSBNF wiring, and the linear appendix proves only the short-yield fact actually needed. Long normalization proofs remain in the appendices, while the former three boundary sections are consolidated into Section 10.

\textbf{Revision location:} Sections 1--2; Sections 4--5; Sections 7--10; Appendices A--C.
\emph{v82 PDF anchors: pp.~1--4, ll.~23--146; pp.~6--10, ll.~244--375; pp.~10--22, ll.~394--814; pp.~22--26, ll.~844--968.}

\subsubsection{Section 1}\label{section-1}

\textbf{R1.3. Comment (quoted excerpt):} ``Stating the paper's place within
the field and the significance of the language families considered would
make the paper more compelling. Conversely, the current text is too
heavily weighted toward technical content and could be made more
concise.''

\textbf{Response:} Agreed. The introduction now states exactly three
contributions: (1) exact finite-witness reconstruction together with the
separate conservative Gold learner; (2) the quantitative fixed-window and
linear characteristic-data results; and (3) the structural position of the
framework relative to the fixed-window hierarchy, natural deterministic CFLs,
and Clark's congruential family. The fixed morphism is presented as a finite
comparison/typing bias, and the fixed-\(h\) learning claim is explicitly
separated from the expressive union RS.

\textbf{Revision location:} Section 1.
\emph{v82 PDF anchor: pp.~1--3, ll.~23--90.}

\textbf{R1.4. Comment (quoted excerpt):} ``The notations \(\mathcal D_L\) and
\(\mathcal C_h^{\mathrm{cf}}\) are not standard and require definitions.
In particular, the phrasing `the general context-free class
\(\mathcal C_h^{\mathrm{cf}}\)' is liable to be misunderstood.''

\textbf{Response:} Agreed. \(\mathcal D_L(x)\) is defined at its first
occurrence in the Introduction and again formally in the preliminaries. The
notation \(\mathcal C_h^{\mathrm{cf}}\) is now defined at its first use in the
Introduction as the class of context-free \(\sim_h\)-substitutable languages;
\(\mathcal C_h^{\mathrm{lin}}\), \(\mathsf{RS}_h\), and \(\mathsf{RS}\) are then
defined formally in the framework section. The phrase ``general context-free
class'' has been removed.

\textbf{Revision location:} Introduction; Section 2.2 and the opening of Section 3.
\emph{v82 PDF anchors: pp.~1--3, ll.~23--90; p.~4, ll.~135--146; pp.~4--6, ll.~147--243.}

\textbf{R1.5. Comment (quoted excerpt):} ``The contributions are enumerated as
(1)-(6), but it is unnecessary to count technical details as separate
contributions. Moreover, at this point the text uses many terms whose
technical meaning has not yet been made clear, and readers may not be
able to follow them.''

\textbf{Response:} Agreed. The list has been reduced to three
substantive contributions and no longer treats proof-local devices as
separate contributions.

\textbf{Revision location:} Section 1.2.
\emph{v82 PDF anchor: pp.~2--3, ll.~69--86.}

\textbf{R1.6. Comment (quoted excerpt):} `` `Yoshinaka's classical language'
is often called `the Łukasiewicz language' in the literature.''

\textbf{Response:} Agreed. We use ``Łukasiewicz language,'' add the
suggested Autebert--Berstel--Boasson (1997) handbook reference, and
state its relation to the one-bracket Dyck language concisely.

\textbf{Revision location:} Section 10.3.
\emph{v82 PDF anchors: pp.~1--3, ll.~23--90; pp.~20--21, ll.~739--777.}

\subsubsection{Section 2}\label{section-2}

\textbf{R1.7. Comment (quoted excerpt):} ``I doubt the claim that `the
learners considered in this paper produce their output independently of
the order in which positive examples are observed.' It appears that the
learning algorithm is being conflated with the grammar-construction
algorithm that it uses as a subroutine. Indeed, the `learning algorithm'
in Section 5 does not converge, and hence does not satisfy Gold's
criterion. Even if the learning algorithm were corrected to follow Clark
\& Eyraud, the output would still differ, when learning, e.g., the
language \(a^n c b^n\), between the case where the characteristic
sample itself is received first and the case where the first positive
example received is very long.''

\textbf{Response:} The reviewer is correct.  The submitted
order-independence claim was false and has been deleted.  The batch operator
and the sequential learner are now formally distinct: only the batch operator
is set-driven.  The stabilized syntactic grammar of the conservative learner
may depend on presentation order, although its generated language eventually
equals \(L\).  The reviewer's example
\(L=\{a^ncb^n:n\ge0\}\) is consistent with exactly this distinction and is
retained in the discussion of Section 4.2.

\textbf{Revision location:} Section 2.1; Section 4.2; Corollary 5.5.
\emph{v82 PDF anchors: pp.~3--4, ll.~117--134; pp.~7--8, ll.~280--286; p.~10, ll.~361--375.}

\textbf{R1.8. Comment (quoted excerpt):} ``This section should thoroughly
define the necessary notation and terminology, including basic notions
such as (linear) context-free grammars - some terms appear to be defined
ad hoc in footnotes. Those terms include `type', `reconstruction',
`realize', `reachable', `observation',
`control', `local configuration', etc. If these words carry no
mathematical definition required to read the paper and were introduced
only to aid intuition, then in the present form they tend to confuse the
reader, and in some cases it would be better not to name them at all. A
flood of excessive terminology can itself burden the reader and impede
smooth understanding.''

\textbf{Response:} Agreed. The preliminaries now define CFGs, the
nonterminal language \(L_G(A)\), syntactic-congruence classes, linear
CFGs, reachable/productive/reduced grammars, useless and nullable
nonterminals, unit productions and unit closure, binarization, terminal
reaching contexts, the fixed empty grammar \(H_0\), text learning,
conservative learning, set-driven reconstruction, characteristic samples
for the batch operator, sample consistency, distributions, and the
fixed-\(h\) classes. Proof-local terminology is introduced only where
used, and the linear section no longer uses wrapper/spine terminology
before its definition. The former ``finite typed reconstruction basis''
terminology has been removed entirely. Several avoidable metavariable
collisions have also been removed where they could impede reading: the
four-element monoid and its concrete homomorphism in Section 9 are now denoted
$M_{\pm,e}$ and $h_{\pm,e}$; the center-symbol metavariable there is $\xi$
rather than the zero-element symbol $\zeta$ used in Section 7; local
fixed-window length parameters in the boundary proofs have descriptive names;
and the intermediate binarized grammar in Appendix A is denoted
$G_{\mathrm{bin}}$.

Section 4.1 now also defines the elements of $\widehat V(K)$ as
``observed nonterminals'' before that phrase is used in proofs.

\textbf{Revision location:} Section 2; Section 4.1; Section 5.1; Section 9.1; Appendix A.2.
\emph{v82 PDF anchors: pp.~3--4, ll.~91--146; pp.~6--7, ll.~245--279; pp.~8--9, ll.~289--330; pp.~17--18, ll.~623--654; pp.~23--24, ll.~873--906.}

\textbf{R1.9. Comment (quoted excerpt):} ``The term `finite typed
reconstruction basis', a combination of undefined words, is clearly
defined here; however, if it denotes an atomic concept for which none of
its four constituent words is individually defined, a shorter term would
be easier to understand.''

\textbf{Response:} Agreed. Rather than rename this proof-local object,
we removed both the term and the basis itself. The revised completeness
argument uses only the yield-typed target grammar and the finite witness
set \(\mathsf W(\widetilde G)\) that is actually needed for exact
reconstruction.

\textbf{Revision location:} Section 5.1.
\emph{v82 PDF anchor: pp.~8--9, ll.~289--330.}

\subsubsection{Section 3}\label{section-3}

\textbf{R1.10. Comment (quoted excerpt):} ``The paper emphasizes (somewhat
excessively and awkwardly) that h is fixed. To match this, RS\_h should
be defined before RS. The family RS is not shown to be learnable from
positive examples (and indeed it is not learnable); what is learnable is
RS\_h.''

\textbf{Response:} Agreed. \(\mathsf{RS}_h\) is defined first and RS is
introduced only for expressive comparisons. The revision now states the
stronger conclusion explicitly: over every nonempty alphabet, neither
\(\mathsf{RS}\) nor \(\mathsf{RS}\cap\mathrm{CFL}\) is identifiable from
positive data, because both contain all regular languages and are therefore
superfinite in Gold's sense.

\textbf{Revision location:} Section 2.2 and the opening of Section 3.
\emph{v82 PDF anchors: p.~4, ll.~135--146; pp.~4--6, ll.~147--243.}

\textbf{R1.11. Comment (quoted excerpt):} ``Definition 3.2 requires care. On
this point, the way Clark \& Eyraud (2007) state their definition is
somewhat careless; however, given that in Section 6.1 of that paper they
state that \(a^+\) is substitutable, and that their grammar-construction
algorithm excludes the empty string, it
is reasonable to interpret the intent of their definition as Yoshinaka's
(2008) 0,0-substitutability.''

\textbf{Response:} We corrected the definition so that fixed-\(h\)
substitutability is stated for nonempty internal fragments \(x,y\in
\Sigma^+\). The empty word is handled separately at the start symbol.
To make this convention explicit at its first use, we have also added a
footnote in the Introduction explaining the point raised by the reviewer:
Clark--Eyraud treat \(a^+\) as substitutable, whereas including
\(\lambda\) among the compared fragments would exclude \(a^+\).
The sufficient-condition discussion in Section 3 now also makes explicit
why the kernel-refinement condition is restricted to \(\Sigma^+\): requiring
\(\ker h\subseteq\equiv_L\) on all of \(\Sigma^*\) would already exclude
\(a^+\) for the trivial monoid. Clark--Eyraud substitutability is then
stated under the same nonempty-fragment convention and identified with the
trivial-monoid case. The revised text now says explicitly that
Clark--Eyraud's Definition~5 is written for all strings, while their
Section~6.1 treats $a^+$ as substitutable; Clark (2013) is cited for the later
explicit nonempty-fragment formulation. In the \((k,\ell)\) definition all
six string variables are explicitly declared to lie in \(\Sigma^*\).

\textbf{Revision location:} Introduction; footnote 1; Definition 2.1 and the Clark--Eyraud special case; the fixed-window definition in Section 3.1.
\emph{v82 PDF anchors: pp.~1--3, ll.~23--90; p.~1, ll.~34--39; p.~4, ll.~139--146; p.~5, ll.~177--190; pp.~5--6, ll.~177--221; p.~5, ll.~186--191.}

\textbf{R1.12. Comment (quoted excerpt):} ``Proposition 3.5. \(\sqcup\) and
\(\mathrm{pref}_k\) are undefined. \ldots{} The proof is very long. If
the multiplication on \(T\) (as a monoid) were given explicitly in
advance, then one line would suffice to explain that \(\tau\) is a
monoid homomorphism. Rephrasing in terms of \(h_{k,\ell}\) or
\(M_{k,\ell}\) is needlessly cumbersome.''

\textbf{Response:} The fixed-window correspondence has been shortened as
suggested. The symbol $\sqcup$ is no longer used. Prefix and suffix notation is
defined before use, and the finite summary monoid is now given with its
multiplication explicitly.
With that multiplication, Proposition 3.3 now defines the homomorphism
$h_{k,\ell}$ directly. A one-bit tag explicitly separates short-word states
from long-word summary states, and the identity
$h_{k,\ell}(xy)=h_{k,\ell}(x)\cdot h_{k,\ell}(y)$ is immediate from
the four product cases. Surjectivity then transfers associativity from
concatenation. The substitution step is stated for $s,t\in\Sigma^+$, matching
the nonempty-fragment convention. The revision also records the converse:
Yoshinaka's $(k,\ell)$ premise factors $py_1q$ and $py_2q$ necessarily have
the same $h_{k,\ell}$-type, so fixed-$h_{k,\ell}$ substitutability yields the
fourth word directly. Thus the correspondence is an equivalence, not merely
an embedding. Immediately afterward we record the unnumbered fixed-window
counterexample criterion used later in Sections 8--9: for sufficiently long
factors, equal length-\(k\) prefixes and length-\(\ell\) suffixes give
equal \(h_{k,\ell}\)-types, so a shared context together with unequal
distributions is already a complete counterexample.

\textbf{Revision location:} Proposition 3.3 and the fixed-window counterexample criterion.
\emph{v82 PDF anchors: pp.~5--6, ll.~192--208; p.~6, ll.~209--221.}

\textbf{R1.13. Comment (quoted excerpt):} ``Rather than saying that \(h\) is
given as data (the meaning is not crystal clear), it would be cleaner to
state early in the paper something like `fix an efficiently computable
\(h:\Sigma^*\to M\)'. Define `linear context-free'.
\ldots{} `Explicit finite monoid homomorphism' is not defined. \ldots{}
Technically unnecessary words should be removed, especially in the
statements of theorems.''

\textbf{Response:} Agreed. We now fix a finite monoid by its
multiplication table and letter images, explicitly state how
\(h\)(\(w\)) is computed, define linear CFGs in the preliminaries, and
use ``fixed efficiently computable homomorphism'' in the main theorem.

\textbf{Revision location:} Sections 2--3.2.
\emph{v82 PDF anchor: pp.~3--6, ll.~91--243.}

\textbf{R1.14. Comment (quoted excerpt):} ``I do not understand `The proof for
fixed \(h\)'. Do you mean `The proof for Theorem 3.7'?''

\textbf{Response:} Agreed. This stray phrase has been deleted. The
revised text now states the relevant theorem and its proof directly,
without this wording.

\textbf{Revision location:} Section 3.2.
\emph{v82 PDF anchor: p.~6, ll.~222--243.}

\subsubsection{Section 4}\label{section-4}

\textbf{R1.15. Comment (quoted excerpt):} ``Definition 4.2. `productive' and
`reachable' are undefined. The phrasing `This section proceeds in the
order typed local control \(\to\) finite typed basis.' is too casual.''

\textbf{Response:} Agreed. These notions are defined in the
preliminaries. The revised typed grammar is explicitly the reduced part
of the full refinement, so thereafter all of its symbols are reachable and
productive by convention and the repeated qualifications are removed. The
quoted informal sentence has also been removed.

\textbf{Revision location:} Section 2; Section 5.1.
\emph{v82 PDF anchors: pp.~3--4, ll.~91--146; pp.~8--9, ll.~289--330.}

\textbf{R1.16. Comment (quoted excerpt):} ``Footnote 2: `productive' is already
used earlier. Define it in Section 2 rather than in a footnote.''

\textbf{Response:} Agreed. ``Reachable'' and ``productive'' are now
defined in the preliminaries, before either term is used in the target
refinement; the former explanatory footnote is gone.

\textbf{Revision location:} Section 2; Section 5.1.
\emph{v82 PDF anchors: pp.~3--4, ll.~91--146; pp.~8--9, ll.~289--330.}

\textbf{R1.17. Comment (quoted excerpt):} ``The text says `the set of reachable and productive typed nonterminals'; is there any nonterminal of $\tilde{G}$ that fails this condition? This may be intended as a helpful reminder or emphasis, but it confuses the reader instead.''

\textbf{Response:} No. The revised notation defines $\widetilde G$ as the reduced part of the full yield-typed refinement, so every symbol retained in $\widetilde G$ is reachable and productive by definition. Repeated qualifications to that effect have been removed.

\textbf{Revision location:} Section 5.1.
\emph{v82 PDF anchor: pp.~8--9, ll.~289--330.}

\textbf{R1.18. Comment (quoted excerpt):} ``The first paragraph states no
reason why `different occurrences must be treated as different typed
copies' at all. Expressions making clear `for what purpose and for what
reason' are needed. Example 4.3 is labeled `(why typing is needed)', yet
why it is needed is not stated. \ldots{} If the authors wish to explain
the reason, the development would be smoother if they first presented
the learning algorithm and then, within its completeness proof, stated:
`here we assume the target grammar is type-refined, and we show that the
learner's hypothesis grammar can mimic it.'\,''

\textbf{Response:} We have adopted the reviewer's suggested order of
presentation. Section 4 first defines the batch reconstruction operator
\(\mathcal B_h\) and the conservative sequential learner
\(\mathcal A_h\), with no target refinement assumed. Section 5 then
introduces the yield-typed target grammar specifically for the
completeness proof and states explicitly that the purpose is to show
that the already-defined hypothesis grammar can mimic typed target
derivations. Outer-context typing has been removed. The one precise role
of the remaining yield typing is then visible inside that simulation:
for a typed binary rule, Rule (R3) must move from the canonical yield
\(\omega(X)\) to \(\omega(Y)\omega(Z)\), and the typed refinement
guarantees \(h(\omega(X))=h(\omega(Y))h(\omega(Z))\). No outer-context
annotation is used.

\textbf{Revision location:} Section 4; Section 5.1 and Theorem 5.3.
\emph{v82 PDF anchors: pp.~6--8, ll.~244--286; pp.~8--9, ll.~289--330; p.~9, ll.~332--351.}

\textbf{R1.19. Comment (quoted excerpt):} ``Lemma 4.4. It is confusing that an
expression concerning derivations in \(G\) appears immediately after `In
the full typed refinement \(\widetilde G^{\mathrm{full}}\),'. It is
unnecessary or should be moved to the end of the sentence.''

\textbf{Response:} The former Lemma 4.4 has been removed in the
reorganization. Its needed content is now folded into Proposition 5.2,
which states language preservation and the yield invariant directly and
keeps derivations of the original grammar and of the typed refinement
syntactically separate.

\textbf{Revision location:} Proposition 5.2.
\emph{v82 PDF anchor: pp.~8--9, ll.~309--331.}

\textbf{R1.20. Comment (quoted excerpt):} ``Lemma 4.8. `If \(X\) is reachable
and productive in \(\widetilde G\),' This likewise confuses the
reader.''

\textbf{Response:} The revised \(\widetilde G\) is defined once as the
reduced part of the full yield refinement, so every symbol of
\(\widetilde G\) is reachable and productive by convention. The former
Lemma 4.8 is deleted, and the redundant hypothesis no longer appears.

\textbf{Revision location:} Section 5.1.
\emph{v82 PDF anchor: pp.~8--9, ll.~289--330.}

\textbf{R1.21. Comment (quoted excerpt):} ``The opening paragraph (the
definitions of \(\omega\) and \(\chi\)) would connect better if placed
immediately before the definition of CS.''

``Since the term `characteristic sample' is defined in Def. 2.1, avoid
writing `Define the characteristic sample by' here unless it is
immediately clear that this set \(CS(\widetilde G)\) satisfies Def.
2.1.''

``The term `realizable' is defined here but used earlier, whereas the
term `realized', used in many places, is undefined.''

``The phrase `canonical observation word' could be defined when
\(CS(\widetilde G)\) is defined.''

\textbf{Response:} Agreed. The canonical yield \(\omega\) and canonical
context \(\chi\) are now named and defined before either notation is used
in the concrete completeness calculation, and the finite witness set is
introduced immediately afterward. To avoid calling the set
characteristic before that property has been proved, it is denoted
\(\mathsf W(\widetilde G)\) and called only a finite witness set at this
stage. Only after Theorem 5.4 establishes exact reconstruction do we
conclude that \(\mathsf W(\widetilde G)\) is a characteristic sample for
the batch operator in the sense defined in the preliminaries. This also
avoids overloading the conventional \(C\)S notation for characteristic
sets/samples. The unused basis, \(NT(\cdot)\), \(\mathrm{Rule}(\cdot)\),
``realized/realizable,'' and ``exhibited'' terminology and the
definition-restating lemmas have been removed.

\textbf{Revision location:} Section 5.1.
\emph{v82 PDF anchor: pp.~8--9, ll.~289--330.}

\textbf{R1.22. Comment (quoted excerpt):} ``Lemma 4.6. This is far too trivial
a corollary. Either label it as a Corollary or delete it.''

``Lemma 4.7. Given what the reader has understood so far, this is
trivial and needs no proof.''

``Theorem 4.9. I cannot understand what this theorem shows.''

``Theorem 4.12. The first sentence should be removed. The second and
third sentences are not understandable owing to the undefined term
`exhibited'.''

\textbf{Response:} The former example and those intermediate results
have been removed. The entire typed-refinement section is now one
definition, one invariant proposition, and the witness-set construction.

\textbf{Revision location:} Section 5.1.
\emph{v82 PDF anchor: pp.~8--9, ll.~289--330.}

\textbf{R1.23. Comment (quoted excerpt):} ``\(\alpha\) is a string, not a
symbol.''

``Said `In particular, if \(X\to YZ\),' but the antecedent of the `if'
is unrelated to its consequent.''

``Lemma 4.5. `For (iii), every derivation in \(G\) erases typed
annotations' I cannot understand this. Do you mean 'For (iii), every
derivation
in \(G\) with typed annotations erased gives \ldots'?''

``Is \(NT(\widetilde G)\) not the same as \(\widetilde V\) on page 9?''

``Definition 4.11. Is \(\mathrm{Rule}(\widetilde G)\) different from
\(\widetilde P\) on page 9?''

\textbf{Response:} Agreed. These sentences and duplicate notations
disappear in the rewrite. The lifting statement is now Proposition 5.2,
stated directly as language preservation plus the yield invariant; the
proof says that erasing type annotations maps a typed derivation to a
derivation of the original grammar. The former trivial corollaries are
deleted.

\textbf{Revision location:} Section 5.1.
\emph{v82 PDF anchor: pp.~8--9, ll.~289--330.}

\subsubsection{Section 5}\label{section-5}

\textbf{R1.24. Comment (quoted excerpt):} ``The learning algorithm as defined
here does not converge. This is the fatal error of the paper.''

\textbf{Response:} The reviewer is correct; this was the central
correctness defect of the submitted learning section. It is corrected by
the constructor/learner separation described above: Section 4.1 defines
the batch constructor, Section 4.2 the conservative sequential learner,
and Corollary 5.5 proves syntactic Gold stabilization.

\textbf{Revision location:} Section 4.2; Corollary 5.5.
\emph{v82 PDF anchors: pp.~7--8, ll.~280--286; p.~10, ll.~361--375.}

\textbf{R1.25. Comment (quoted excerpt):} ``Corollary 5.8. This cannot be
accepted as a proof. By the definition in Section 5.1, the learner's
hypothesis grammar keeps growing as it receives more positive
examples.''

\textbf{Response:} This objection arises from the same defect, so we
changed the learner rather than only the proof. The revised learner is
conservative and Corollary 5.5 proves that the hypothesis grammar itself
eventually stabilizes syntactically.

\textbf{Revision location:} Section 4.2; Corollary 5.5.
\emph{v82 PDF anchors: pp.~7--8, ll.~280--286; p.~10, ll.~361--375.}

\textbf{R1.26. Comment (quoted excerpt):} ``Remark 5.1. I cannot grasp the
point. What does `soundness of a rule' mean? \ldots{} If one sentence is
to be added, something like `the grammar is designed so that the
language derived by each nonterminal \([x:u,v]\) should
be the set of strings substitutable with \(x\)' might aid intuition.
Stating a mathematical property concisely helps the reader more than
deploying adjectives that carry no mathematical definition.''

``Lemmas 5.3 and 5.4 merely repeat definitions and should be deleted.''

\textbf{Response:} The remark and definition-restating lemmas are
deleted. In place of the old informal remark, immediately after Rules
(R1)--(R5) we now give the simple intended reading suggested by the
reviewer: \([x:u,v]\) represents possible replacements
for the observed factor \(x\) in context (\(u\), \(v\)); the following
soundness invariant makes precise that derived replacements retain
the same \(h\)-type and distributional behavior. The rule labels were
changed to (R1)--(R5) so that references such as ``Rule (R3)'' cannot be
confused with numbered displayed equations. The revised order is now:
batch grammar, sample consistency and soundness, sequential learner,
then target-side yield refinement and completeness.

\textbf{Revision location:} Section 4 (Rules (R1)--(R5) and the soundness
invariant); Section 5.2.
\emph{v82 PDF anchors: pp.~6--8, ll.~244--286; pp.~9--10, ll.~331--375.}

\textbf{R1.27. Comment (quoted excerpt):} ``Lemma 5.2. In the second sentence,
the first claim is trivial and the second is not understandable. What is
the corresponding local configuration? Its correspondence with the part
following `More precisely' is unclear.''

\textbf{Response:} Agreed. The former Lemma 5.2 and the undefined
``corresponding local configuration'' terminology have been removed.
Their role is replaced by the explicit rule-by-rule soundness invariant
in Theorem 4.2, which states directly what every hypothesis nonterminal
may derive and proves it by induction on derivation-tree height, with a case analysis on the root reconstruction rule.

\textbf{Revision location:} Theorem 4.2.
\emph{v82 PDF anchor: p.~7, ll.~261--279.}

\textbf{R1.28. Comment (quoted excerpt):} ``Page 14, line 1: `Since \(G\) is in
SSBNF' \(\to\widetilde G\)?''

\textbf{Response:} Corrected during the rewrite; the relevant grammar is
named explicitly throughout the completeness proof.

\textbf{Revision location:} Theorem 5.3.
\emph{v82 PDF anchor: p.~9, ll.~332--351.}

\textbf{R1.29. Comment (quoted excerpt):} ``Theorem 5.7. Make the statement
mathematical. What do you mean by `a finite sample exposes the finite
typed reconstruction basis'? A sample is a string set and the FTRB is a
tuple.''

\textbf{Response:} Theorem 5.4 now states the exact implication
\(\mathsf W(\widetilde G)\subseteq K\subseteq L\Rightarrow L(\mathcal B_h(K))=L\)
with all assumptions in the theorem statement.

\textbf{Revision location:} Theorem 5.4.
\emph{v82 PDF anchor: pp.~9--10, ll.~352--360.}

\subsubsection{Section 6}\label{section-6}

\textbf{R1.30. Comment (quoted excerpt):} ``If the paper adopts the stance of
Remark 6.4, then the fact that the grammar construction can be carried
out in polynomial time is almost trivial, so the content of this section
could be presented in a single paragraph.''

\textbf{Response:} Section 6 is now a single counting proof with
\(n_K=\sum_{w\in K}(|w|+1)\): \(O(n_K^2)\) observed nonterminals,
\(O(n_K^3)\) factorization rules, at most \(O(n_K^4)\) candidate
transport/type pairs, and a safe explicit-output bound \(O(n_K^5)\) when
string keys are copied. This size convention also counts \(\lambda\)
correctly.

\textbf{Revision location:} Section 6.
\emph{v82 PDF anchor: p.~10, ll.~376--393.}

\subsubsection{Section 7}\label{section-7}

\textbf{R1.31. Comment (quoted excerpt):} ``This is a special case of the
arguments in the preceding sections and could be much shorter \ldots{}
Too long explanation may blur the focus --- what the reader should
understand, and what differs from the preceding sections.''

``Lemma 7.3. `strict linear derivation' is not defined.''

\textbf{Response:} We substantially shortened the material corresponding
to submitted Section 7; it is now Section 9. As explained in the
Overview, only the statements of Proposition 9.1 and Lemma 9.2
remain in the main text, because Theorem 9.3 uses their exact
single-spine form and the resulting short-yield bound directly; their
detailed proofs are now in Appendices B and C. The undefined ``strict
linear derivation'' terminology has been removed, and Proposition 9.1
now states explicitly that the binary rule contains one fresh
terminal-wrapper child. At the end of Section 6, a compact
doubling-grammar example explains why no grammar-size-only polynomial
characteristic-data bound is claimed for arbitrary CFG representations.

\textbf{Revision location:} The general-CFG limitation in Section 6; Proposition 9.1 and the surrounding discussion; Lemma 9.2; Appendix B; Appendix C.
\emph{v82 PDF anchors: p.~10, ll.~376--393; p.~16, ll.~604--610; p.~17, ll.~611--612; pp.~24--25, ll.~907--945; pp.~25--26, ll.~946--968.}

\textbf{R1.32. Comment (quoted excerpt):} ``Proposition 7.2 does not need two
full pages on the proof. I do not know whether exactly the same
restriction has been defined before, but it is a combination of existing
restrictions that are known to be inessential, and most readers will
accept Proposition 7.2 as trivial or as folklore.''

``Lemmas 7.5 and 7.6 are too trivial.''

\textbf{Response:} The former two-page normalization proof and the
elementary intermediate lemmas have been removed from the main
development. Proposition 9.1 and Lemma 9.2 remain only as the two
statements used immediately by Theorem 9.3; their detailed proofs are in
Appendices B and C, respectively. The former elementary lemmas are
absorbed into those appendix proofs or deleted.

\textbf{Revision location:} Proposition 9.1 and the surrounding discussion; Lemma 9.2; Theorem 9.3 and the following clarification; Appendix B; Appendix C.
\emph{v82 PDF anchors: p.~16, ll.~604--610; p.~17, ll.~611--612; p.~17, ll.~613--622; pp.~24--25, ll.~907--945; pp.~25--26, ll.~946--968.}

\subsubsection{Sections 8--10}\label{sections-810}

\textbf{R1.33. Comment (quoted excerpt):} ``The choice of the alphabet
\(\{d,u,;\}\) is unfortunate. The semicolon in particular is very hard
to read within the running text, and \(u\) is a source of confusion
because it almost always denotes a string elsewhere in the paper.''

``The proof of Theorem 8.6 should be simplified. Introducing complex
formulas and lavishing Greek letters in an attempt to shorten the
witness is not the essence of the theorem and only burdens the reader.''

\textbf{Response:} Agreed. The counter alphabet is now
\(\{\uparrow,\downarrow,\#\}\) and the cap is \(\rho\). We also isolate
in Section 3 a short fixed-window counterexample criterion: two sufficiently
long factors with the same length-\(k\) prefix and length-\(\ell\) suffix,
a shared context, and different distributions already refute
\((k,\ell)\)-substitutability by Proposition 3.3. Section 10.2 then uses only
the two words \((\uparrow\downarrow)^k\uparrow^{\rho-1}\#^\ell\) and
\((\uparrow\downarrow)^k\uparrow^\rho\#^\ell\): they share the empty
context, while one extra leading \(\uparrow\) distinguishes their
distributions. No six-variable replacement witness is expanded there.  The revised
manuscript also thanks the anonymous reviewer immediately after the theorem
for prompting this direct two-word presentation.

\textbf{Revision location:} Section 10.2.
\emph{v82 PDF anchors: pp.~1--3, ll.~23--90; pp.~19--20, ll.~714--738.}

\textbf{R1.34. Comment (quoted excerpt):} ``There is no need to devote an
entire subsection to this material. It is unclear what Corollary 8.7
establishes. \ldots{} Proposition 3.5 and Theorem 8.6 are meaningful;
but folding them into Corollary 8.7 actually obscures the paper's
important contributions.''

`` `We note that a block of a word in CCL refers to each of the subwords
delimited by the separator symbol ;.' This is a new definition, not a
note. In fact, the proof of Theorem 8.6 uses `block' of a CCL string in
an entirely different sense.''

``Sections 8-10 all concern the richness of the class of the learning
target languages; consolidating them into a single section and keeping
the explanations and proofs concise would help the reader grasp the
overall picture.''

\textbf{Response:} We followed the consolidation suggestion. The former standalone \(\Delta^*\)
development has been condensed to a single proposition, retaining only the nonlinear witness
needed for the scope claim and for retaining the submitted nonlinear
separation from the fixed-window hierarchy. The former three sections
are now one boundary section.  The former Corollary 8.7 has been removed as a
numbered result.  After the capped-counter theorem the revision records only
the fixed-slice consequence
\(\mathrm{CTR}_\rho\in\mathcal C_{h_\rho}^{\mathrm{cf}}
\setminus\mathsf{KL}\) for a finite recognizing homomorphism $h_\rho$;
the target-dependent union $\mathsf{RS}$ is retained only as an expressive
umbrella, not as a learnable class.  The nonlinear-witness proof now uses
``segment'' rather than reusing the reviewer-flagged term ``block''. The
capped-counter family has been renamed from \(\mathrm{CCL}_p\) to
\(\mathrm{CTR}_\rho\), with the alphabet changed from \(\{d,u,;\}\) to
\(\{\uparrow,\downarrow,\#\}\) for readability. The capped-counter
regularity proof is one paragraph, the uncapped-counter and Dyck
exclusions follow from a single finite-monoid obstruction lemma, and the
Łukasiewicz relation is stated by citation rather than reproved. Instead, a compact proposition records exactly the two boundary facts
needed here: \(\Delta^*\), with \(\Delta:=\{a^n b^n:n\ge0\}\), is a
nonlinear context-free member of a fixed-\(h\) class and is not
\((k,\ell)\)-substitutable for any fixed \(k,\ell\). Its fixed-\(h\)
proof is organized by a single admissible-entry-height set rather than repeated
casewise height calculations, and its fixed-window exclusion invokes the same
counterexample criterion from Section 3. Thus \(L_{\pm,e}\) carries the
nonregular linear separation from the fixed-window hierarchy, while
\(\Delta^*\) retains the nonlinear separation. The final summary now states the alphabets explicitly: the regular separation
is over $\{\uparrow,\downarrow,\#\}$, the nonregular linear separation is
over $\{a,b,c,d,e\}$, and the nonlinear separation is over $\{a,b\}$.
The two optional appendices have been removed entirely.

\textbf{Revision location:} Section 10.
\emph{v82 PDF anchors: pp.~1--3, ll.~23--90; pp.~18--22, ll.~655--814.}

\textbf{R1.35. Comment (quoted excerpt):} ``Definition 9.1. Define the symbol
\#, and make its use consistent in Section 9.''

``Proposition 9.2 is trivial and requires no proof. Proposition 9.6 is
also a well-known fact, as already stated in Ginsburg \& Greibach,
Deterministic context-free languages, Information and Control 9 (1966),
620-648.''

``Lemma 9.8. The relationship between the Łukasiewicz language and the
Dyck language is also well known.''

``The proof of Proposition 10.3 may also be omitted or greatly
shortened. In the proof of Proposition 10.5, consider aligning the
symbols with those of Definition 3.3.''

\textbf{Response:} Agreed. The separator is now uniformly \# and its
counter action is defined once. Standard facts are either stated briefly
or cited; in particular, the one-bracket Dyck language is now cited
directly to Ginsburg--Greibach (1966) when its deterministic
context-free status is stated, the Łukasiewicz language is also cited
there for deterministic-CFL background, and Autebert--Berstel--Boasson
(1997) is cited for the standard Dyck/Łukasiewicz relationship. The
submitted Lemma 9.8 on that standard relationship is therefore no longer
reproduced as a standalone lemma. The former boundary arguments are compressed into the finite-monoid
obstruction, the short fixed-word quotient lemma, and the compact nonlinear
\(\Delta^*\) proposition described above. The repeated fixed-window
six-variable witnesses are replaced by the single counterexample criterion of
Section 3, and the \(\Delta^*\) proof uses one admissible-entry-height
invariant. Notation has been aligned with the definitions; the former
standalone \(\Delta^*\) development remains condensed into that single
proposition.

\textbf{Revision location:} Section 10.
\emph{v82 PDF anchors: pp.~1--3, ll.~23--90; pp.~18--22, ll.~655--814.}

\textbf{R1.36. Comment (quoted excerpt):} ``Lemma 9.11 is a fundamental
property of \(\sim_h\) substitutability; since the basic mathematical
properties of \(\sim_h\) substitutability are collected in an appendix,
placing it there is one option.''

\textbf{Response:} We followed the reviewer's suggestion to isolate the
fundamental property, but did not move the reduced statement to an appendix.
We retained only the right-quotient form actually used and placed it
immediately before its single application. The submitted Lemma 9.11 has been reduced to the
right-quotient case actually used by the Łukasiewicz argument; this reduced
statement is now Lemma 10.7 and appears immediately before that application.
The submitted boundary-property appendices were removed in the consolidation
of Sections 8--10; the three proof appendices are devoted to the general
thickness-preserving normalization and the two detailed linear proofs. Keeping the
two-line quotient argument next to its sole use makes the dependency
local without recreating a separate boundary appendix.

\textbf{Revision location:} Lemma 10.7; proof appendices.
\emph{v82 PDF anchors: p.~21, ll.~769--777; pp.~22--26, ll.~844--968.}

\subsubsection{Section 11}\label{section-11}

\textbf{R1.37. Comment (quoted excerpt):} ``It is worth mentioning that
\(\mathsf{RS}\cap\mathrm{CFL}\) is a subset of the class of
Clark's congruential CFLs.''

\textbf{Response:} Added as Proposition 10.8. Clark (2010) already observes that substitutable and \((k,\ell)\)-substitutable languages lie in his congruential family and uses the Dyck language as a standard example beyond ordinary substitutability; the revision now says this explicitly. Our proposition extends the containment statement to every fixed-\(h\) slice and hence to \(\mathsf{RS}\cap\mathrm{CFL}\). The revised statement is made
slice-wise: for every fixed finite-monoid homomorphism \(h\),
\(\mathcal C_h^{\mathrm{cf}}\subseteq\mathsf{CONG}(\Gamma)\); taking the union
over \(h\) gives \(\mathsf{RS}(\Gamma)\cap\mathrm{CFL}
\subseteq\mathsf{CONG}(\Gamma)\). For a nonempty target in a fixed
\(h\)-slice, each non-start yield-typed nonterminal generates nonempty strings
of one \(h\)-type and shares a reaching context, so fixed-\(h\)
substitutability places all of its yields in a single syntactic-congruence
class. Removing the separated start symbol and using its finitely many typed
successors as the initial set therefore gives a congruential grammar; the
notation \(L_G[I]\) makes this finite-initial-set convention explicit. The
inclusion is proper over \(\{a,b\}\) because \(D_1\) is congruential but
lies outside RS.

\textbf{Revision location:} Proposition 10.8.
\emph{v82 PDF anchor: pp.~21--22, ll.~788--814.}

\subsection{Reviewer 2}\label{reviewer-2}

\subsubsection{\texorpdfstring{General positioning of
fixed-\(h\)}{General positioning of fixed-h}}\label{general-positioning-of-fixed-h}

\textbf{R2.1. Comment (quoted excerpt):} ``At a general level, the paper
offers an abstraction of the learning of Yoshinaka's
\((k,\ell)\)-substitutable languages by replacing the fixed-length
window by an arbitrary recognizable congruence, and lifts the known
results accordingly. The abstraction places the problem in a powerful
and well-understood framework, but it requires a whole morphism to be
known in advance for each language, whereas the earlier settings fixed a
single uniform parameter valid across the entire class. The paper
acknowledges that learning the morphism is outside its scope, yet the
interest of the setting -- beyond providing a more abstract framework
for the study of \((k,\ell)\)-substitutable languages -- remains
unclear. A promising direction would be to deepen the links with
pushdown automata, which Sections 8 to 10 already touch on. The links
with the work of Takada \ldots{} are also worth investigating. In a more
practical direction, this work could be read as extending automaton
typing \ldots{} to context-free substitutable languages, and the
conditions a typing must satisfy could be studied.''

\textbf{Response:} We agree that the modeling role of fixed
$h$ needed to be stated more sharply.  The revision now fixes one homomorphism
$h$ uniformly over the whole class $\mathcal C_h^{\mathrm{cf}}$, just as one
fixes a single pair $(k,\ell)$ in Yoshinaka's setting; selecting or learning
$h$ is a separate problem and remains outside the present scope.  Following
the reviewer's suggestion, we describe $h$ as a finite comparison/typing bias
and make the comparison with Coste et al. (2004) more precise.  Their expert
typing is encoded by finite-state type automata: the type of an occurrence may
depend on its left and right contexts, and state merging is constrained so
that the typing embedded in the inferred automaton remains compatible with
the expert-supplied typing.  Our object is formally different---$h$ gives a
compositional finite type to terminal factors rather than an occurrence-level
contextual type to automaton structure---so the revised manuscript presents
the relation as a methodological analogy, not an identification of the two
models.

The revision also makes both the gain and the limitation of the abstraction
concrete.  Proposition 3.3 identifies the fixed $h_{k,\ell}$ slice exactly
with Yoshinaka's $(k,\ell)$ class.  Section 9 gives a four-element fixed
typing for the nonregular linear language $L_{\pm,e}$, which lies outside
every fixed-window class.  Conversely, Section 7 shows that arbitrary fixed
typings do not in general inherit the ordinary-thickness polynomial
characteristic-data bound: Theorem 7.7 gives the set-driven lower bound.
The tree representation, residual height, clean subtrees, shortcut nodes, and
antichains used in that proof are now defined explicitly before the
substitutability lemma.

We also revised the control-set comparison after checking the complete
sources.  Takada's 1988 result reduces even-linear inference to regular
control-set inference on a fixed universal even-linear grammar.  The 1995
chapter makes this architecture especially explicit: on a fixed universal
unambiguous even-linear grammar, each target has a unique canonical regular
control set, and learning even-linear languages reduces to learning regular
languages.  The same chapter develops regular-control representations more
broadly for linear grammars and equal-matrix grammars.  These approaches use
a target-specific regular language of production labels to constrain
derivations; fixed $h$, by contrast, is a class-level finite bias on terminal
factors and constrains only which factors may be compared.  We therefore
claim neither equivalence nor subsumption nor a class separation with respect
to the control-set methodology in general.  Koshiba--M\"akinen--Takada
(1997) is cited separately for the later positive-data restriction.

Finally, Section 9 records the classical one-turn-PDA characterization of
linear languages, while Section 10 isolates a finite-monoid obstruction that
excludes the uncapped counter language and $D_1$ from every finite-typing
slice.  These additions are intended as concrete connections, not as a full
pushdown characterization; that broader problem remains open.

\textbf{Revision location:} Introduction; Proposition 3.3; Section 7.2;
Section 9; Sections 10.2--10.4; Conclusion.
\emph{The structural locators above are current.  Because the Introduction
was refined again after the v82 PDF was compiled, the page/printed-line
anchors for this response item will be regenerated from the final revised PDF
before submission.}

\subsubsection{Main concern 1: outer context typing and
transport}\label{main-concern-1-outer-context-typing-and-transport}

\textbf{R2.2. Comment (quoted excerpt):} ``My main concern is the typed
refinement, and specifically transport. Rule (2) allows a nonterminal to
move between arbitrary contexts with no condition whatever on their
types. Is this intended? If the outer components $m,n$ record where a
nonterminal may appear, why does the rule that relocates it ignore them
entirely? As far as I can tell the components are then read by no proof:
Lemma 4.6 is invoked only for the yield component, and replacing
$A^{m,n}_{p}$ by $A_{p}$ throughout leaves every result intact. Does the
author agree, and if not, where exactly is an outer type consulted? The
situation seems to me to follow from Definition 2.3 itself, whose algebraic
condition bears on the substituted fragment and imposes nothing on the
surrounding context -- so that a learner mirroring that definition has
nothing to guard transport with. If that is right, then the typed transport
announced in Section 1.1 does not occur anywhere in the article, and the
claim that occurrences must be separated according to their outer types
cannot be established by any example. The yield typing, which is genuinely
essential and is what becomes vacuous under the trivial monoid, should not
be left in a footnote.''

\textbf{Response:} The reviewer is correct. The submitted outer-context annotations were not used by any proof, so the submitted claim that such outer typing was necessary was unsupported. We have therefore removed those annotations and withdrawn the description of the mechanism as ``typed transport,'' together with the claim that occurrences must be separated according to their outer types. Rule (R2) transports the same observed factor $x$ only between contexts in which $x$ itself has been observed; that common occurrence supplies the shared-context premise needed for substitutability, and no $h$-type condition on the surrounding contexts is used. The target refinement now contains only yield-typed symbols \(A_\mu\). This agrees with the reviewer's
observation that the algebraic condition in the submitted Definition 2.3
constrains the substituted fragment, not the surrounding context. The hypothesis
grammar still records observed contexts syntactically as \([x:u,v]\). The
genuinely necessary yield typing, previously left in a footnote, is now introduced
explicitly in the main text of Section 5.1 and is used at one precise completeness
step. For a typed binary rule
\(X=A_{\mu\nu}\to B_\mu C_\nu\), Rule (R3) changes the canonical
representative \(\omega(X)\) to \(\omega(B_\mu)\omega(C_\nu)\), and the
typed refinement supplies the required equality of \(h\)-values. Thus
Section 4 defines the reconstruction rules first, and Section 5
introduces exactly the target typing needed to simulate them.

\textbf{Revision location:} Section 4; Section 5.1 and Theorem 5.3.
\emph{v82 PDF anchors: pp.~6--8, ll.~244--286; pp.~8--9, ll.~289--330; p.~9, ll.~332--351.}

\subsubsection{Main concern 2: extent of the linear
theorem}\label{main-concern-2-extent-of-the-linear-theorem}

\textbf{R2.3. Comment (quoted excerpt):} ``The linear result raises a
different question: how large is the class it covers? The three closing
sections situate RS, not \(\mathcal C_h^{\mathrm{lin}}\), so the reader
is given no sense of what Theorem 7.9 reaches beyond the linear
substitutable languages. The capped-counter languages are regular and
hence
linear, so the separation from Yoshinaka's hierarchy already holds
inside a linear class, whereas \(\{a^n b^n\}^*\) is not linear, so the
final section establishes nothing at that level. As it stands the linear
theorem reads as a technical strengthening on a class of unknown extent,
the non-regular linear case being left open.''

\textbf{Response:} We thank the reviewer for identifying this missing
case. Section 9 now adds a nonregular linear language \(L_{\pm,e}\) with a
four-element monoid morphism. Proposition 9.4 proves that it is linear,
nonregular, fixed-\(h\) substitutable, and neither ordinarily substitutable nor
\((k,\ell)\)-substitutable for any fixed \(k,\ell\). Thus Theorem 9.3 has
nonregular linear content beyond the fixed-window hierarchy. Its quantitative
proof has also been shortened: the linear normal form yields
\(\tau_h^{\mathrm{typ}}\le n_t\), after which the general typed-thickness
corollary supplies the characteristic-data bound. The fixed-window
quantitative comparison is handled separately in Section 8. New Section 7
explains the general fixed-\(h\) quantitative boundary in terms of typed
thickness, gives a fixed two-element source-to-typed exponential gap, and then
proves the stronger set-driven ordinary-thickness lower bound of Theorem 7.7
using nested finite targets with compact reduced CFGs. The general-CFG
limitation remains documented by the doubling-grammar example in Section 6,
which now explicitly notes that each singleton target in that family is
\(\sim_h\)-substitutable for every \(h\). Correspondingly, we no longer
describe the corrected sequential learner as identifying the linear class in
``polynomial time and data.'' The revised linear theorem states polynomial
characteristic data for the batch reconstruction, while Gold identification is
supplied separately by Corollary 5.5.

For completeness, Section 10.1 retains the compact nonlinear witness
\(\Delta^*\), showing that the separation also persists at a nonlinear
context-free target.

\textbf{Revision location:} The general-CFG limitation in Sections 6--7; Section 8; Theorem 9.3 and Proposition 9.4; Proposition 10.1.
\emph{v82 PDF anchors: pp.~10--15, ll.~376--560; pp.~15--16, ll.~561--599; p.~17, ll.~613--622; pp.~17--18, ll.~635--654; pp.~18--19, ll.~680--707.}

\subsubsection{Further comments}\label{further-comments}

\textbf{R2.4. Comment (quoted excerpt):} ``The paper would gain considerably
from being simplified: much of Section 4 is definitional, several
results restate what their proofs have just established, and objects are
regularly used before they are introduced, which forces the reader to
hold material in suspense. Relatedly, since $\tilde G$ is by
construction the trimmed refinement of $\tilde G_{\mathrm{full}}$,
its symbols are reachable and productive by definition; repeating this
qualification at every occurrence adds length without adding content,
and one sentence recording the convention would suffice.''

\textbf{Response:} The repeated qualifications are removed and the
reduced-part convention is stated once.

\textbf{Revision location:} Section 5.1.
\emph{v82 PDF anchor: pp.~8--9, ll.~289--330.}

\textbf{R2.5. Comment (quoted excerpt):} ``Proposition 3.1: the hypothesis
\(L=h^{-1}(F)\) characterises the regular languages exactly. Replace `In
particular' by `Thus' (and the proof could make explicit why).''

\textbf{Response:} Agreed. Proposition 3.1 now states the finite-monoid
characterization directly and concludes: ``Thus every regular language
belongs to RS.'' The proof cites the standard finite-monoid characterization
and then makes explicit the consequence relevant here: if
\(L=h^{-1}(\mathsf{Acc})\) and \(h(x)=h(y)\), then
\(h(uxv)=h(uyv)\) for every \(u,v\), so the two distributions agree.

\textbf{Revision location:} Proposition 3.1.
\emph{v82 PDF anchor: p.~5, ll.~171--176.}

\textbf{R2.6. Comment (quoted excerpt):} ``Definition 4.11: the basis
\(B(\widetilde G)\) is never used. Is it useful to introduce it as well
as Theorem 4.12 and Remark 4.13? (Notations NT() and Rule() are
unnecessary since it is sufficient to define \(G\) = () and use defined
\(V\) and \(P\).)''

\textbf{Response:} Agreed. These objects and notations have been
removed.

\textbf{Revision location:} Section 5.1.
\emph{v82 PDF anchor: pp.~8--9, ll.~289--330.}

\textbf{R2.7. Comment (quoted excerpt):} ``Example 4.3 is not in SSBNF''

\textbf{Response:} The former example has been removed entirely; the
revised proof no longer needs it.

\textbf{Revision location:} Section 5.1.
\emph{v82 PDF anchor: pp.~8--9, ll.~289--330.}

\textbf{R2.8. Comment (quoted excerpt):} ``Lemma 5.2 (i): If \(X\) carries
several terminal rules \(\omega(X)\) is not sufficient, and rule (3)
should be used (proof needs to be corrected accordingly).''

\textbf{Response:} The reviewer is correct; this was a genuine gap in
the submitted completeness argument. The witness set now contains
$u_Xav_X$ for every typed terminal rule $X\to a$, ensuring that the observed
nonterminal $[a:u_X,v_X]$ is present. The terminal base case then applies
Rule (R3) when the canonical yield differs from $a$, followed by Rule (R4).

\textbf{Revision location:} Theorem 5.3.
\emph{v82 PDF anchor: p.~9, ll.~332--351.}

\textbf{R2.9. Comment (quoted excerpt):} ``Are Lemmas 5.3 and 5.4 needed?''

\textbf{Response:} Agreed. They are deleted.

\textbf{Revision location:} Sections 4--5.1.
\emph{v82 PDF anchors: pp.~6--9, ll.~244--330; pp.~6--10, ll.~244--375.}

\textbf{R2.10. Comment (quoted excerpt):} ``Corollary 5.8 Should begin with
`For a fixed homomorphism \(h\) from \(\Sigma^*\) to a finite monoid
\(M\), \ldots.', or `Being given a homomorphism \(h\) from\ldots.'\,''

\textbf{Response:} The revised Gold-identification statement now begins
exactly in the requested explicit form: ``For a fixed homomorphism \(h\)
: \(\Sigma^*\to M\) into a finite monoid \(M\), \ldots{}''.

\textbf{Revision location:} Corollary 5.5.
\emph{v82 PDF anchor: p.~10, ll.~361--375.}

\textbf{R2.11. Comment (quoted excerpt):} ``section 7.1 Transformation to
standard SSLNF (3 pages) could be condensed or given in appendix.''

\textbf{Response:} We have taken the appendix option suggested by the
reviewer. As summarized in the Overview, Section 9 retains only the statements
needed directly by Theorem 9.3, while the detailed normalization and
short-yield proof are moved to Appendices B and C. The main text therefore
contains only the statements and obtains the polynomial characteristic-data
bound from the general typed-thickness corollary.

\textbf{Revision location:} Proposition 9.1 and the surrounding discussion; Lemma 9.2; Theorem 9.3 and the following clarification; Appendix B; Appendix C.
\emph{v82 PDF anchors: p.~16, ll.~604--610; p.~17, ll.~611--612; p.~17, ll.~613--622; pp.~24--25, ll.~907--945; pp.~25--26, ll.~946--968.}

\medskip
\noindent\textbf{Reproducibility note.}
A Lean 4 formalization accompanying parts of the development is archived on
Zenodo (DOI 10.5281/zenodo.22939434). Its scope predates and does not include
the set-driven ordinary-thickness lower-bound results of Section 7. The
manuscript does not rely on the formalization as a substitute for any paper
proof.

\end{document}
